\chapter{Background}
\label{ch:background}

\section{Conformal Prediction}

The theoretical foundation of CP~\cite{vovk2005algorithmic}
rests on the assumption of exchangeability of the data points \((X_i,Y_i)\).
Exchangeability, a weaker condition than independence and identical distribution (i.i.d.),
requires that the joint distribution of the data remains invariant under permutation.
Under this assumption, CP leverages the idea of conformity to assess
how well a new, unseen data point fits with previously observed data.
A nonconformity score \(s(x, y) \in \mathbb{R}\),
quantifies how ``atypical'' a potential prediction is relative to past examples.

The choice of the nonconformity measure (NCM) is critical
for the practical utility of CP~\cite{kato2023measures}.
While the validity guarantee holds for any NCM,
an informative NCM leads to smaller, more useful prediction sets.
Common NCMs often leverage the output of the underlying model.
For classification, examples include
the inverse predicted probability of the true class \((1 - \hat \mathbb{P}(y | x))\)
or the margin between the probability of the true class
and the highest probability of any other class.

The full CP procedure is computationally intensive,
requiring recalculation of scores (and potentially model retraining)
for every test point and every possible label.
To overcome this, the split conformal prediction
(also known as Inductive Conformal Prediction or ICP)
method was developed~\cite{papadopoulos2002inductive} and is now the most popular variant.
Its main drawback is reduced data efficiency, as a portion of the data
must be held out for calibration and is not used for model training.
ICP works as follows:

\begin{enumerate}
\item Split the available data into a training set \(\mathcal{D}_{train}\) and a
  calibration set \(\mathcal{D}_{cal}\). The training set is used to train the model,
  while the calibration set is used to compute the nonconformity scores.

\item Train the underlying prediction model using only \(\mathcal{D}_{train}\).

\item Compute the nonconformity scores \(s_i = s(X_i, Y_i)\) for all examples
  \((X_i, Y_i)\) in the calibration set \(\mathcal{D}_{cal}\) using the model
  trained in step 2.

\item Determine the threshold \(\hat q\) by calculating the \(\frac{\lceil
  (n_{cal} + 1)(1 - \alpha) \rceil}{n_{cal}}\) empirical quantile of the calibration
  scores \(\{s_i\}_{i \in \mathcal{D}_{cal}}\), where \(n_{cal} = |\mathcal{D}_{cal}|\).

\item For a new test input \(X_{test}\), calculate the nonconformity score
  \(s(X_{test}, y)\) for each possible label \(y\) using the fixed model. The
  prediction set is then all labels where the nonconformity score is less than or
  equal to the threshold: \(\mathcal{C}(X_{test}) = \{y \in Y : s(X_{test}, y) \leq \hat q\}\).
\end{enumerate}


\subsection{Evaluating Sets}

Evaluating the performance and utility of prediction sets generated by Conformal
Prediction (CP) involves assessing several key aspects beyond just the
underlying model's predictive accuracy:

\paragraph{Validity (Coverage)}
The primary theoretical guarantee provided by CP is marginal coverage. Under the
exchangeability assumption, for any chosen significance level \(\alpha \in (0, 1)\) the
prediction set \(\mathcal{C}(X_{n+1})\) generated by the conformal algorithm is guaranteed to
contain the true label \(Y_{n+1}\) with probability at least \(1 - \alpha\):

\begin{equation}
  \label{eq:coverage}
   \mathbb{P}(Y_{n+1} \in \mathcal{C}(X_{n+1})) \geq 1 - \alpha.
\end{equation}

This is typically evaluated empirically by
calculating the percentage of test samples for which
the true label falls within the predicted set.
This empirical coverage should ideally match or exceed the nominal coverage level \(1 - \alpha\).

% While standard CP guarantees marginal coverage, research also
% explores achieving conditional coverage, ensuring validity across specific
% subgroups or conditional on certain features, though this is often more
% challenging. (Need citation for this)

\paragraph{Efficiency (Set Size)}
While guaranteed coverage is essential,
the practical usefulness of prediction sets hinges on their efficiency,
typically measured by their size.
Smaller sets provide more precise information and are generally preferred.
Evaluation usually involves calculating the average set size across a test dataset.
Some evaluations might also consider the informativeness of sets, for instance,
by analysing the frequency of overly large or ``trivial'' sets
(e.g., sets containing almost all possible labels).

\paragraph{Adaptiveness}
Ideally, prediction sets should be adaptive,
meaning their size reflects the uncertainty associated with a specific instance.
The set should be smaller for instances where the model is confident
and potentially larger for ambiguous instances where uncertainty is higher.
Evaluating adaptiveness involves analysing the relationship
between set size and measures of instance difficulty or model confidence.

\subsection{CP in Medical Imaging}

While standard Conformal Prediction (CP) provides foundational guarantees,
refinements have been developed to enhance its practicality for complex classification tasks.
Methods like Adaptive Prediction Sets (APS)~\cite{romano2020aps} specifically address
the limitation that raw classifier outputs often yield miscalibrated probabilities,
leading to prediction sets that don't achieve the target coverage level.
APS corrects this by using a calibration set to determine
the appropriate score threshold needed to empirically satisfy the desired coverage.

However, even with guaranteed coverage,
APS can sometimes produce overly large prediction sets,
particularly in domains like image classification
where models might assign small, noisy probabilities to many incorrect classes.
Addressing this, Regularized Adaptive Prediction Sets (RAPS)~\cite{angelopoulos2020raps}
introduced a regularization penalty for unlikely classes
(those ranked lower in the probability list).
This regularization effectively reduces the influence of noisy tail probabilities,
resulting in significantly smaller and more stable prediction sets
without sacrificing the formal coverage guarantees.

Despite these advancements in achieving reliable marginal coverage
\((\mathbb{P}(Y \in \mathcal{C}(X)))\)
and optimizing set size, methods like APS and RAPS fundamentally target
the overall error rate of \emph{miscoverage}.
They do not, by default, offer direct control over
specific, potentially more critical types of errors,
such as the rate of false negatives for high-risk conditions.
This control over specific, asymmetric risks
is precisely the challenge addressed by Conformal Risk Control (CRC).

While these methods offer general advancements for CP in image tasks,
specific adaptations tailored to the nuances of medical imaging have also emerged.
For instance, Lu et al.\@ explored the use of CP in dermatology image classification,
focusing on achieving `fairness' in medical AI systems~\cite{lu2022fair}.
Their work adapts standard CP methods to provide equalized coverage across skin types on a
skin lesion classification task, ensuring that the prediction sets contain the
true label with the desired probability (e.g., 90\%) regardless of the subgroup
membership.  This addresses important ethical considerations regarding
algorithmic bias but primarily focuses on equitable coverage guarantees across
groups rather than controlling specific types of prediction errors for the
overall population or for high-risk conditions.

Another application tackles the challenge of ordinal classification in medical imaging,
specifically for grading the severity of spinal stenosis from MRI scans~\cite{lu2022improving}.
They developed an adaptation of APS to generate prediction sets that are
contiguous intervals of severity levels.
This ensures that the true severity grade is captured
within the predicted range with high probability,
which is more intuitive for ordinal tasks than potentially disjoint sets from standard CP.
However, like fair CP, this approach focuses on ensuring
coverage of the true label (or rank) within the set.

Applications of CP in medical imaging have primarily focused on ensuring that
the true label is contained within the prediction set with a certain frequency.
They do not directly offer mechanisms to control the rate of
specific, clinically critical errors, such as
minimizing the False Negative Rate (FNR) for detecting a severe condition.
This is the gap that Conformal Risk Control (CRC) aims to fill,
by controlling the expected value of a user-defined, clinically relevant loss function,
moving beyond general coverage guarantees to manage specific diagnostic risks.

% \cite{lu2022fair} skin lesions, three dermatological categories - non-neoplastic, benign, malignant.
% \cite{lu2022improving} spinal stenosis in lumbar spine MRI, four severities - none, mild, moderate, severe.

\section{Conformal Risk Control}
